% Source presentations for register row N-c-023.

\subsection{Stable models as candidate-dependent reducts}
\label{sec:p09-presentation-n-c-023}

\subsubsection{The Gelfond--Lifschitz construction by complete calculation}

A finite ground normal rule has the form
\[
 H\leftarrow A_1,\ldots,A_m,
   \mathsf{not}\ B_1,\ldots,\mathsf{not}\ B_n,
\]
where all displayed expressions other than \(\mathsf{not}\) are ground atoms.
Let \(\operatorname{At}(P)\) be the finite set of atoms occurring in a finite
ground program \(P\), and let \(M\subseteq\operatorname{At}(P)\) be a candidate
interpretation.  Unlike the full Herbrand base, \(\operatorname{At}(P)\) is only
this syntactically occurring finite set; atoms outside it cannot be generated by
a reduct and hence cannot belong to a stable model.  The reduct \(P^M\) is
formed rule by rule:
\begin{enumerate}
\item delete the rule if \(B_j\in M\) for at least one \(j\);
\item otherwise retain its head and positive body and erase all
  \(\mathsf{not}\ B_j\).
\end{enumerate}
The result is a positive definite program.  Therefore it has the unique least
model
\[
 \Gamma_P(M)=\operatorname{lfp}(T_{P^M}),
\]
and Definition~\ref{def:p09-stable-model} says exactly
\[
 M\text{ is stable}
 \quad\Longleftrightarrow\quad
 M=\Gamma_P(M).
\]
Thus the least fixed point is taken only \emph{after} the candidate-dependent
deletion step.  It is not an iteration of one monotone operator that directly
assigns a model to the original normal program.

The map reverses inclusion.  If \(M\subseteq N\), every rule retained in
\(P^N\) is retained in \(P^M\), so \(P^N\subseteq P^M\) and
\(\Gamma_P(N)\subseteq\Gamma_P(M)\).  This antimonotonicity is why the adjacent
well-founded construction uses the monotone composite \(\Gamma_P^2\); it does
not give \(\Gamma_P\) itself a least fixed point.

For a nonground rule, the standing of the construction is equally explicit:
first take its declared ground instances and only then form the reduct.  If
function symbols generate infinitely many ground terms, the finite decision
arguments below no longer apply.  If the program is function-free and its
constants are finite, its grounding is finite.

For the program
\[
 P=\{p\leftarrow\mathsf{not}\ q,\quad
     q\leftarrow\mathsf{not}\ p\},
\]
the candidate dependence is visible by checking all four candidates:
\[
\begin{array}{c|c|c}
M&P^M&\operatorname{lfp}(T_{P^M})\\
\hline
\varnothing&\{p\leftarrow,\ q\leftarrow\}&\{p,q\}\\
\{p\}&\{p\leftarrow\}&\{p\}\\
\{q\}&\{q\leftarrow\}&\{q\}\\
\{p,q\}&\varnothing&\varnothing
\end{array}
\]
Hence the stable models are exactly \(\{p\}\) and \(\{q\}\).  In contrast,
\[
 P'=\{p\leftarrow\mathsf{not}\ p\}
\]
has none: the empty candidate produces the fact \(p\), while the candidate
\(\{p\}\) deletes the only rule and produces the empty least model.  These
calculations establish both nonuniqueness and nonexistence without importing a
solver.

For a positive definite program \(Q\), the reduct is \(Q\) for every candidate,
so \(\Gamma_Q(M)=M_Q\) for every \(M\).  Consequently
\[
 M\text{ is a stable model of }Q
 \quad\Longleftrightarrow\quad M=M_Q.
\]
This is the precise adjacent case connecting stable semantics to
Definition~\ref{def:p09-definite-program}; it does not turn a program with
default negation into a monotone fixed-point construction.

Every stable model is an ordinary two-valued model of its program and, more
strongly, a minimal Herbrand model.  This is Theorem~1 of the original paper.
For ordinary modelhood, a rule deleted from the reduct has a false
default-negated body condition in the candidate, while a retained rule is
satisfied because the candidate models its positive reduct.  For minimality,
suppose that a proper subset \(N\subsetneq M\) were also a Herbrand model of
the original program.  Every rule retained in \(P^M\) came from a rule whose
default-negated atoms are outside \(M\), hence outside \(N\).  Since \(N\)
models that original rule, it models the retained positive rule.  Thus \(N\)
would be a model of \(P^M\), contradicting that \(M\) is the least model of
\(P^M\) \cite[\S2, thm.~1]{GelfondLifschitz1988}.

The converse fails, and stable semantics is not Clark completion.  For
\[
 R=\{p\leftarrow p\},
\]
every reduct is \(R\), whose least model is \(\varnothing\); hence
\(\varnothing\) is the unique stable model.  The completion is instead the
tautology \(p\leftrightarrow p\), which has both classical interpretations
\(\varnothing\) and \(\{p\}\).  The extra classical model \(\{p\}\) is
self-supported rather than generated by the positive least-model construction
\cite{GelfondLifschitz1988}.

The live first-party publication page still identifies the 1988 paper and
points to Lifschitz's former author-hosted PostScript file.  That dead link has
an archived capture containing the complete eleven-page author copy
corresponding to published pp.~1070--1080.  All eleven pages, including the
proofs supplied for
Theorems~1 and~3 and the intervening lemma, have now been inspected.
Section~2 first calls the fixed points above \emph{stable
sets}; after Theorem~1 proves minimality, it licenses the name \emph{stable
models}.  The paper then reserves the formulation \emph{the stable model
semantics is defined for \(P\)} for the case in which \(P\) has exactly one
stable model, which it declares canonical.  The broader modern convention used
in this handbook speaks of the stable-model relation even when a program has
zero or several stable models.  The reduct is unchanged; only the later
terminology has widened.  Lifschitz's retrospective independently reproduces
that reduct as Definition~C, including grounding before the two deletion steps
\cite[\S5, p.~43]{Lifschitz2008TwelveDefinitions}.

\paragraph{Comparisons proved or recorded in the original paper.}
Theorem~2, explicitly credited there to Van Gelder, Ross, and Schlipf, states
in the paper's then-current terminology that if \(P\) has a well-founded model,
that model is its unique stable model.  Corollary~1 specializes this to locally
stratified programs and their perfect models; Corollary~2 specializes further
to stratified programs and their iterated-fixed-point models.  The handbook's
three-valued formulation is
Theorem~\ref{thm:p09-well-founded-properties}: its uniqueness conclusion is
the total case \(L=G\), not an identification of every partial well-founded
interpretation with a two-valued stable model
\cite[\S4, thm.~2 and cors.~1--2]{GelfondLifschitz1988}.

The remaining comparison uses the stable expansions of
Definition~\ref{def:p08-autoepistemic-expansion}.  For a ground program \(P\),
let \(\mathcal I(P)\) be the autoepistemic theory obtained by reading each rule
as an implication and replacing each default-negated body atom
\(\mathsf{not}\ B\) by \(\neg L B\), where \(L\) is the belief operator.
This is a metalinguistic translation into another logic, not a second reduct.
Original Theorem~3 says that if \(P\) has the unique stable model \(M\), then
\(\mathcal I(P)\) has a unique stable expansion \(E\) and
\[
 M=E\cap\operatorname{At}(P).
\]
Its proof reduces the claim to the lemma that \(E\) is a stable expansion of
\(\mathcal I(P)\) exactly when it is a stable expansion of the objective
reduct \(P^{E\cap\operatorname{At}(P)}\).  Objective theories have their unique
expansion determined by ordinary propositional closure, so any competing
expansion would make its objective atoms a competing stable model of \(P\)
\cite[\S4, thm.~3]{GelfondLifschitz1988}.  None of these 1988
results is the later finite-program complexity classification below.

\subsubsection{One fixed function-free NP representation}

Let \(\Pi\) be a search problem over finite instances, let \(I\) range over its
instances, and let \(S_\Pi(I)\) be its finite set of solutions.  Marek and Truszczyński's
uniform presentation separates three objects:
\begin{enumerate}
\item a varying extensional database \(\operatorname{edb}_{\Pi}(I)\) over one
  fixed relational schema;
\item one finite function-free program \(P_\Pi\), fixed for every instance of
  \(\Pi\);
\item a polynomially computable map \(\operatorname{sol}^{I}_{\Pi}\) from
  stable models of \(\operatorname{edb}_{\Pi}(I)\cup P_\Pi\) to solutions in
  \(S_\Pi(I)\), with the one-to-one condition stated in their definition.
\end{enumerate}
The database is the input, the program is the invariant problem description,
and the last map decodes rather than changes a model.  These kinds cannot be
collapsed: a stable model contains the input facts and chosen auxiliary atoms,
whereas a decoded solution is the finite object those choices represent.
Their Theorem 5.2 says that a decision problem has such a uniform SLP
presentation exactly when it is in \(\mathsf{NP}\)
\cite[\S5, thm.~5.2]{MarekTruszczynski1999}.

Three-colorability gives a complete small instance of the pattern.  Let the
varying database contain \(node(v)\) for every vertex and \(edge(u,v)\) for
every represented edge.  The fixed color-choice rules are
\[
\begin{aligned}
 red(X)&\leftarrow node(X),\mathsf{not}\ green(X),
                         \mathsf{not}\ blue(X),\\
 green(X)&\leftarrow node(X),\mathsf{not}\ red(X),
                           \mathsf{not}\ blue(X),\\
 blue(X)&\leftarrow node(X),\mathsf{not}\ red(X),
                          \mathsf{not}\ green(X).
\end{aligned}
\]
For each color predicate \(c\), add the fixed conflict rule
\[
 bad\leftarrow edge(X,Y),c(X),c(Y),\mathsf{not}\ bad.
\]
The atom \(bad\) has no other defining rule.  Its use is necessary because the
normal syntax in force has no empty or disjunctive head: this odd negative cycle
eliminates exactly the candidates whose positive body violates the constraint.

Take the graph with vertices \(u,v\) and the one edge \((u,v)\).  Write
\(E_G\) for its three database facts and consider
\[
 M_{rg}=E_G\cup\{red(u),green(v)\}.
\]
In the reduct relative to \(M_{rg}\), the retained choice rules for the two
vertices are
\[
 red(u)\leftarrow node(u),
 \qquad
 green(v)\leftarrow node(v).
\]
Every other ground color-choice rule is deleted.  The conflict rules remain
positive because \(bad\notin M_{rg}\), but none has a true same-color edge
body.  Forward closure therefore gives exactly \(M_{rg}\), so this candidate is
stable.  By contrast, the candidate
\[
 M_{rr}=E_G\cup\{red(u),red(v)\}
\]
retains a positive rule deriving \(bad\), and hence is not the least model of
its reduct.  Adding \(bad\) to the candidate does not repair it: then every
conflict rule is deleted, leaving \(bad\) underived.

The calculation generalizes in both directions.
\begin{description}
\item[From a stable model to a coloring.]
Every database fact occurs in every reduct and hence in every stable model.
For each vertex, a candidate containing no color retains all three choice
rules, whose reduct derives all three colors.  A candidate containing two or
three colors deletes every rule capable of deriving any of those simultaneous
choices.  Stability is therefore possible only with exactly one color per
vertex.  No stable candidate contains \(bad\), because a candidate containing
it deletes its only defining rules.  If the chosen colors agree across an
edge, the retained conflict rule derives \(bad\), again contradicting
stability.  The selected colors consequently form a proper coloring.
\item[From a coloring to a stable model.]
Given a proper coloring, take the database facts together with its one color
atom per vertex and omit \(bad\).  The reduct retains exactly the rule deriving
each selected color.  Every conflict-rule body is false, so positive closure
adds nothing else.  The candidate is the least model of its reduct.
\end{description}
The two constructions are inverse: once the input facts are fixed, a stable
model is determined by its color atoms.  Decoding those atoms is therefore a
bijection between stable models and proper three-colorings.

The six displayed rules are fixed; only the \(node\) and \(edge\) facts grow.
After grounding on a graph with \(n\) vertices, there are three choice rules per
vertex and at most three conflict rules per ordered vertex pair, so the ground
program has polynomial size.  Since graph three-colorability is
\(\mathsf{NP}\)-complete, this is also a direct polynomial reduction to stable-
model existence for finite normal programs.

The general theorem has the same two proof roles at a more abstract level.  In
the only-if direction, a fixed function-free program has a constant bound on
the number of variables in each rule.  Grounding it over the active input
constants is therefore polynomial in the database size.  A nondeterministic
algorithm guesses a candidate, forms its positive reduct, computes the least
model by forward closure, and applies the polynomial decoding map.  In the
if direction, the program fixed for \(\Pi\) represents the choices in an
accepting polynomially bounded computation, while the varying database carries
the particular input and the required finite position and successor facts;
fixed rules enforce the local transition and acceptance conditions.  The full
uniform compilation is the theorem cited above; the coloring program is a
complete instance of its generate-and-eliminate pattern.  The paper leaves the
corresponding characterization of arbitrary search problems open, and no part
of the theorem says that all stable models can be enumerated in polynomial time.

\subsubsection{Answer-set syntax and complexity boundaries}

The modern ASP overview extends the original language in two independent ways.
First, an explicitly or strongly negated literal \(\neg a\) asserts a negative
fact; it is not the default test \(\mathsf{not}\ a\).  In the common ground
presentation, \(a\) and \(\neg a\) are distinct signed literals and an answer
set is required to be consistent, so it does not contain both.  Second, a
disjunctive rule may have the form
\[
 a_1\lor\cdots\lor a_k
 \leftarrow b_1,\ldots,b_m,
   \mathsf{not}\ c_1,\ldots,\mathsf{not}\ c_n.
\]
Its reduct uses the same two deletion steps, but the result is a positive
\emph{disjunctive} program.  A candidate is an answer set exactly when it is an
inclusion-minimal model of that reduct.  When \(k=1\), the positive reduct is
definite and its inclusion-minimal model is its unique least model, recovering
the stable-model definition above.  These are language-specific model
conditions, not alternate names for every ordinary model
\cite[pp.~93--97]{BrewkaEiterTruszczynski2011}.

For finite propositional normal programs, the three complexity statements in
Theorem~\ref{thm:p09-stable-complexity} can be separated by short reductions.
Existence is in \(\mathsf{NP}\): guess \(M\), form \(P^M\), compute its least
model by positive forward closure, and compare the result with \(M\).  The
three-coloring construction above proves \(\mathsf{NP}\)-hardness.

Credulous consequence asks whether some stable model contains the queried
atom.  Membership uses the same guessed candidate.  For hardness, given \(P\),
add a fresh fact \(w\leftarrow\).  The stable models of the enlarged program
are exactly \(M\cup\{w\}\) for stable models \(M\) of \(P\).  Hence \(P\) has a
stable model exactly when \(w\) is credulously entailed by the enlarged
program.

Skeptical consequence asks whether every stable model contains the queried
atom.  Its complement has the polynomial witness consisting of a stable model
that omits the atom, so the problem is in \(\mathsf{coNP}\).  For hardness, use
the ordinary universal convention that a statement is true in every member of
an empty class.  Given \(P\), add a fresh positive loop
\(w\leftarrow w\).  This makes \(w\) part of the program but never puts it in a
least positive reduct model.  Thus the enlarged program skeptically entails
\(w\) exactly when \(P\) has no stable model.  This reduces the complement of
stable-model existence and proves \(\mathsf{coNP}\)-hardness
\cite{DantsinEiterGottlobVoronkov2001}.

\begin{table}[htbp]
\centering\scriptsize
\begin{tabularx}{\textwidth}{@{}l X X@{}}
\toprule
Input language & Decision or size boundary & Qualification \\
\midrule
finite propositional normal program & \(\mathsf{NP}\)-complete &
candidate reduct has one polynomially computable least model \\
finite propositional disjunctive program &
\(\Sigma^{\mathsf P}_2\)-complete &
candidate reduct may have many minimal models \\
nonground predicate program & grounding may be exponential &
if the rule set itself varies, its number of variables need not be bounded;
ground size and source size are then different measures \\
fixed function-free program plus varying finite database & polynomial grounding
in data size & every fixed rule has a constant number of variables; this is the
uniform setting of Theorem 5.2 \\
\bottomrule
\end{tabularx}
\caption{Syntax-indexed answer-set decision and grounding boundaries.}
\label{tab:p09-asp-complexity}
\end{table}

The overview supplies the distinctions and membership intuition but is not the
primary proof of each hardness theorem.  For a positive disjunctive reduct,
checking that a candidate has no smaller model is a coNP task; existentially
choosing the candidate accounts for the second-level membership intuition.
The \(\Sigma^{\mathsf P}_2\)-hardness result remains cited rather than disguised
as a locally reconstructed quantified reduction.  Solver brands and
architectures lie outside this logical presentation.
